%! Suppress = MultipleIncludes
\documentclass[tikz,crop]{standalone}

\onlyifstandalone{
    \usepackage[outputdir=../../out/]{minted}
    \setminted{autogobble=true}
}

% Part of the preamble, for TikZ figures.
% This is used in both the main document and in the subfigures.
% One exception is minted: since the path depends on the file, it is not set.
\usepackage{tikz}
\usepackage{xcolor}
\usepackage{pgfplots}
\usepackage{tikz-uml}

\pgfplotsset{compat=1.18}
\usepgfplotslibrary{statistics}

\usetikzlibrary{shapes,arrows,positioning,backgrounds,calc,intersections,calc,svg.path}

\definecolor{ugent-re}{RGB}{220, 78, 40}        % vermilion			/ vermiljoen
\definecolor{ugent-we}{RGB}{45, 140, 168}       % no match
\definecolor{ugent-ge}{RGB}{232, 94, 113}       % rose				/ bleekrood
\definecolor{ugent-ea}{RGB}{111, 113, 185}      % distant blue		/ verblauw
\definecolor{ugent-pp}{RGB}{251, 126, 58}       % deep orange		/ dieporanje
\definecolor{ugent-ps}{RGB}{113, 168, 96}       % yellow green		/ geelgroen

\tikzstyle{python}=[fill=ugent-ps!50!white]
\tikzstyle{java}=[fill=ugent-we!50!white]
\tikzstyle{haskell}=[fill=ugent-ea!50!white]
\tikzstyle{js}=[fill=ugent-pp!50!white]
\tikzstyle{c}=[fill=ugent-re!50!white]

\newlength{\block}
\setlength{\block}{0.75cm}

\tikzstyle{a}=[anchor=north west]
\tikzstyle{box}=[a,draw,rectangle]
\tikzstyle{node}=[a,draw,minimum height=0.5cm,align=center,fill=white,text depth=.25ex]
\tikzstyle{document}=[node,tape,tape bend top=none]
\tikzstyle{cont}=[box,minimum height=1\block,minimum width=1\block]
\tikzstyle{arrow}=[draw, -latex]
\tikzstyle{inner}=[box,draw=gray]

% Blue box style
\tikzstyle{bluebox}=[draw=ugent-we,java]
\tikzstyle{redbox}=[draw=ugent-re,c]
\tikzstyle{greenbox}=[draw=ugent-ps,python]

% Some things specific to TESTed imagery.
\tikzstyle{tc}=[box,draw=ugent-ps]
\tikzstyle{comp}=[box,draw=ugent-re,fill=ugent-re,fill opacity=0.05]
\tikzstyle{exec}=[box,draw=ugent-we,fill=ugent-we,fill opacity=0.10]

% Stuff from tested-engine/concept.tex
\tikzstyle{process}=[node,rectangle]
\tikzstyle{terminator}=[node,rectangle,rounded corners=0.5cm]
\tikzstyle{io}=[node,trapezium,trapezium left angle=70,trapezium right angle=-70,minimum width=2.5cm,trapezium stretches=true]
\tikzstyle{small}=[font=\footnotesize,color=darkgray]
\tikzstyle{submission}=[document,align=right,minimum width=3cm,minimum height=1cm,text depth=0.5cm,inner sep=0.5mm,font=\scriptsize]

% Stuff from chatper3/flow.tex
\tikzstyle{height}=[minimum height=0.75\block]
\tikzstyle{contt}=[cont,minimum height=0.75\block]
\tikzstyle{compop}=[comp,text opacity=1]
\tikzstyle{execop}=[exec,text opacity=1]

\tikzstyle{hnode}=[draw,anchor=center,minimum height=\block,text depth=.25ex,align=center]
\tikzstyle{executable}=[hnode,ultra thick,fill=gray!10]
\tikzstyle{inner-exec}=[node,anchor=center,minimum width=3.25\block,densely dotted,font=\footnotesize,fill=none]
\tikzstyle{stmt}=[node,anchor=center,fill=gray!30,minimum width=4.5\block,font=\footnotesize]
\tikzstyle{fieldset}=[minimum height=\block,fill=white,text depth=.5ex,fill=white]

\definecolor{c9ed9e6}{RGB}{158,217,230}
% TODO: create better version of this, or use PDF directly...
\newcommand*\drawtesttube{%
    \kern0.2em\relax
    \begin{tikzpicture}[scale=0.5,baseline=3pt,rotate=100]
%        \path[fill=c9ed9e6] (0.926, -0.5821)arc(229.9677:328.957:0.1058).. controls (1.1377, -0.5027) and (1.1113, -0.4498) .. (1.0583, -0.3969) -- (0.5821, -0.0265)arc(49.3987:229.3987:0.122) -- (0.4498, -0.2117) -- (0.0794, -0.7144)arc(143.1301:233.1301:0.2646) -- (0.1852, -1.1113)arc(232.7082:275.401:0.2646).. controls (0.4498, -1.1642) and (0.5027, -1.1113) .. (0.5556, -1.0583) -- (0.926, -0.5821) -- cycle(0.3704, -0.926) -- (0.3175, -0.926) -- (0.2646, -0.8996) -- (0.2646, -0.8467) -- (0.635, -0.3704) -- (0.7408, -0.4498) -- (0.3704, -0.926) -- cycle;
        \draw[fill=c9ed9e6,draw=none] svg {M35.043,22.196c0.813,0.616,1.77,0.915,2.717,0.915c1.358,0,2.703-0.615,3.589-1.783c1.501-1.98,1.111-4.803-0.868-6.305
			L21.868,0.915c-1.981-1.503-4.803-1.113-6.304,0.867s-1.113,4.804,0.868,6.305l0.254,0.192L2.776,26.633
			c-1.623,2.144-2.315,4.788-1.948,7.45c0.367,2.663,1.749,5.021,3.889,6.643l2.068,1.568c1.767,1.34,3.877,2.045,6.057,2.045
			c0.462,0,0.927-0.031,1.392-0.096c2.663-0.365,5.023-1.747,6.647-3.89L34.79,22.001L35.043,22.196z M13.709,34.918
			c-0.22,0.29-0.502,0.383-0.701,0.41c-0.201,0.026-0.496,0.014-0.785-0.205l-2.07-1.57c-0.289-0.218-0.381-0.5-0.409-0.697
			s-0.014-0.494,0.205-0.783l13.91-18.355l3.759,2.85L13.709,34.918z};
    \end{tikzpicture}%
    \kern0.2em\relax\space
}

% Minted environments for use in Tikz
\newminted[tikzjava]{java}{autogobble,linenos=false,fontsize=\tiny,stripall}
\newminted[tikzpython]{python}{autogobble,linenos=false,fontsize=\tiny,stripall}
\newminted[tikztext]{text}{autogobble,linenos=false,fontsize=\tiny,stripall}


\begin{document}

% Define a tikzcode environment for code.
\newminted[tikzcode]{java}{autogobble,linenos=false,fontsize=\tiny,stripall}
\newminted[tikztcode]{text}{autogobble,linenos=false,fontsize=\tiny,stripall}

\begin{tikzpicture}[x=0.75cm,y=0.75cm]
    % \draw[step=1,gray!10,thin] (0,0) grid (14,-20);
    % Ensure width is the same
%    \node[box,opacity=0,minimum width=13.5\unit] at (0,-1.75) (ts1) {};

    \begin{scope} [shift={(1.875,1.5)}]
        \node[cont] at (0.75, -1.75) (ts1c1) {C\textsubscript{1}};
        \node[cont] at (2.25, -1.75) (ts1c2) {C\textsubscript{2}};
        \node[cont] at (3.75, -1.75) (ts1c3) {C\textsubscript{3}};
        \node[tc,minimum height=2\unit,minimum width=4.5\unit] at (0.5,-1) (tc1) {};
        \node[below right,text=ugent-ps] at (tc1.north west) {\small Unit 1};

        \node[cont] at (5.50, -1.75) (ts1c4) {C\textsubscript{4}};
        \node[cont] at (7.00, -1.75) (ts1c5) {C\textsubscript{5}};
        \node[tc,minimum height=2\unit,minimum width=3\unit] at (5.25,-1) (tc2) {};
        \node[below right,text=ugent-ps] at (tc2.north west) {\small Unit 2};

        \node[cont] at (8.75, -1.75) (ts1c6) {C\textsubscript{6}};
        \node[cont] at (10.25, -1.75) (ts1c7) {C\textsubscript{7}};
        \node[cont] at (11.75, -1.75) (ts1c8) {C\textsubscript{8}};
        \node[tc,minimum height=2\unit,minimum width=4.5\unit] at (8.5,-1) (tc3) {};
        \node[below right,text=ugent-ps] at (tc3.north west) {\small Unit 3};

        \node[comp,minimum height=2.5\unit,minimum width=13\unit] at (0.25,-0.75) (comp) {};

        \node[exec,minimum height=2.25\unit,minimum width=4.25\unit] at (0.625,-0.875) (execu1) {};
        \node[exec,minimum height=2.25\unit,minimum width=1.25\unit] at (5.375,-0.875) (execu2) {};
        \node[exec,minimum height=2.25\unit,minimum width=6\unit] at (6.875,-0.875) (execu3) {};
    \end{scope}

    \begin{scope} [shift={(0,-3)}]
        \node[document,minimum width=17.25\unit, minimum height=8.5\unit,draw=ugent-re] at (0,0) (executable) {};
        \node[below right,text=ugent-re] at (executable.north west) {\small Executable};

        \node[box,draw=ugent-re,minimum height=3.25\unit,minimum width=4\unit] at (0.25,-0.75) (main) {};
        \node[below right,text=ugent-re] at (main.north west) (mtitle) {\small \texttt{main} function};
        \node[below right,text width=3.75\unit] at (mtitle.south west) {%
            \fvset{listparameters=\setlength{\topsep}{0pt}\setlength{\partopsep}{0pt}}%
            \begin{tikzcode}
                switch(argument) {
                  case "1" -> exec1();
                  case "2" -> exec2();
                  case "3" -> exec3();
                  default -> throw ...;
                }
            \end{tikzcode}
        };

        \node[exec,minimum height=3.25\unit,minimum width=4\unit] at (4.5,-0.75) (exec1) {};
        \node[below right,text=ugent-we] at (exec1.north west) (e1title) {\small \texttt{exec1()} function};
        \node[below right,text width=3.75\unit] at (e1title.south west) {%
            \fvset{listparameters=\setlength{\topsep}{0pt}\setlength{\partopsep}{0pt}}%
            \begin{tikzcode}
                outputContextSecret();
                context1();
                outputContextSecret();
                context2();
                outputContextSecret();
                context3();
            \end{tikzcode}
        };

        \node[exec,minimum height=3.25\unit,minimum width=4\unit] at (8.75,-0.75) (exec2) {};
        \node[below right,text=ugent-we] at (exec2.north west) (e2title) {\small \texttt{exec2()} function};
        \node[below right,text width=3.75\unit] at (e2title.south west) {%
            \fvset{listparameters=\setlength{\topsep}{0pt}\setlength{\partopsep}{0pt}}%
            \begin{tikzcode}
                outputContextSecret();
                context4();
            \end{tikzcode}
        };

        \node[exec,minimum height=3.25\unit,minimum width=4\unit] at (13,-0.75) (exec3) {};
        \node[below right,text=ugent-we] at (exec3.north west) (e3title) {\small \texttt{exec3()} function};
        \node[below right,text width=3.75\unit] at (e3title.south west) {%
            \fvset{listparameters=\setlength{\topsep}{0pt}\setlength{\partopsep}{0pt}}%
            \begin{tikzcode}
                outputContextSecret();
                context5();
                outputContextSecret();
                context6();
                outputContextSecret();
                context7();
                outputContextSecret();
                context8();
            \end{tikzcode}
        };

        % Normally, context functions should be here, however, they are drawn at the end.
    \end{scope}


    \draw[-latex,draw=ugent-we] (execu1) to [out=-90,in=90] (exec1);
    \draw[-latex,draw=ugent-we] (execu2) to [out=-90,in=90] (exec2);
    \draw[-latex,draw=ugent-we] (execu3) to [out=-90,in=90] (exec3);

%    \draw[-latex] (ts1c1) to [out=-90,in=90] (cf1);
%    \draw[-latex] (ts1c2) to [out=-90,in=90] (cf2);
%    \draw[-latex] (ts1c8) to [out=-90,in=90] (cf8);

    \draw[-latex,draw=ugent-re] (comp) -- (executable) node[midway,fill=white,fill opacity=0.8,text opacity=1] {\small Code generation and compilation};

    \begin{scope} [shift={(0,-13)}]
        \node[document,draw=ugent-we,minimum width=5\unit, minimum height=3.25\unit] at (0.625, 0) (res1) {};
        \node[below right,color=ugent-we] at (res1.north west) (res1title) {\small \texttt{output1.txt}};
        \node[below right,text width=3.75\unit] at (res1title.south west) {%
            \fvset{listparameters=\setlength{\topsep}{0pt}\setlength{\partopsep}{0pt}}%
            \begin{tikztcode}
                context_separator (C1)
                testcase_separator
                result1
                testcase_separator
                result2
            \end{tikztcode}
        };

        \node[document,draw=ugent-we,minimum width=5\unit, minimum height=3.25\unit] at (6.125, 0) (res2) {};
        \node[below right,color=ugent-we] at (res2.north west) (res2title) {\small \texttt{output2.txt}};
        \node[below right,text width=3.75\unit] at (res2title.south west) {%
            \fvset{listparameters=\setlength{\topsep}{0pt}\setlength{\partopsep}{0pt}}%
            \begin{tikztcode}
                context_separator (C1)
                testcase_separator
                result1
                context_separator (C2)
                testcase_separator
                result1
                ...
            \end{tikztcode}
        };

        \node[document,draw=ugent-we,minimum width=5\unit, minimum height=3.25\unit] at (11.625, 0) (res3) {};
        \node[below right,color=ugent-we] at (res3.north west) (res3title) {\small \texttt{output3.txt}};
        \node[below right,text width=3.75\unit] at (res3title.south west) {%
            \fvset{listparameters=\setlength{\topsep}{0pt}\setlength{\partopsep}{0pt}}%
            \begin{tikztcode}
                ...
                context_separator (C8)
                testcase_separator
                result1
            \end{tikztcode}
        };
    \end{scope}

    \draw[-latex,draw=ugent-we] (exec1) to [out=-90,in=90] (res1);
    \draw[-latex,draw=ugent-we] (exec2) to [out=-90,in=90] (res2);
    \draw[-latex,draw=ugent-we] (exec3) to [out=-90,in=90] (res3);

    % Context functions, to be before the arrows
    \begin{scope} [shift={(0,-3)}]
        \node[cont,minimum width=4.5\unit, minimum height=3\unit,fill=white] at (0.5, -4.75) (cf1) {};
        \node[below right] at (cf1.north west) (c1title) {\small \texttt{context1()} function};
        \node[below right,text width=3.75\unit] at (c1title.south west) {%
            \fvset{listparameters=\setlength{\topsep}{0pt}\setlength{\partopsep}{0pt}}%
            \begin{tikzcode}
                outputTestcaseSeparator();
                testcase1();
                outputTestcaseSeparator();
                testcase2();
                outputTestcaseSeparator();
                testcase3();
                ...
            \end{tikzcode}
        };

        \node[cont,minimum width=4.5\unit, minimum height=3\unit,fill=white] at (5.5, -4.75) (cf2) {};
        \node[below right] at (cf2.north west) (c2title) {\small \texttt{context2()} function};
        \node[below right,text width=3.75\unit] at (c2title.south west) {%
            \fvset{listparameters=\setlength{\topsep}{0pt}\setlength{\partopsep}{0pt}}%
            \begin{tikzcode}
                outputTestcaseSeparator();
                testcase1();
                outputTestcaseSeparator();
                testcase2();
            \end{tikzcode}
        };

        \node[cont,draw opacity=0,align=center] at (10.625, -4.75) {\ldots};

        \node[cont,minimum width=4.5\unit, minimum height=3\unit,fill=white] at (12.25, -4.75) (cf8) {};
        \node[below right] at (cf8.north west) (c8title) {\small \texttt{context8()} function};
        \node[below right,text width=3.75\unit] at (c8title.south west) {%
            \fvset{listparameters=\setlength{\topsep}{0pt}\setlength{\partopsep}{0pt}}%
            \begin{tikzcode}
                outputTestcaseSeparator();
                testcase1();
            \end{tikzcode}
        };
    \end{scope}

    \draw[draw opacity=0] (executable) -- (res2) node[midway,fill=white,fill opacity=0.8,text opacity=1] {\small Execution};

\end{tikzpicture}

\end{document}
